\documentclass[11pt]{article}
\usepackage[margin=1.15in]{geometry}
\usepackage{amsmath,amssymb,amsthm}
\usepackage{booktabs}
\usepackage{xcolor}
\usepackage[colorlinks=true,linkcolor=blue!60!black,citecolor=blue!60!black,urlcolor=blue!60!black]{hyperref}
\hypersetup{pdftitle={Adaptive Planar Search},pdfauthor={Peter Cotton}}

\newtheoremstyle{itheads}{\topsep}{\topsep}{}{}{\itshape}{.}{.5em}{}
\theoremstyle{itheads}
\newtheorem{theorem}{Theorem}
\newtheorem{proposition}{Proposition}

\newcommand{\E}{\mathbb{E}}
\newcommand{\R}{\mathbb{R}}

\title{Adaptive Planar Search:\\[4pt]
{\large A Terminal-Placement Acquisition for Derivative-Free Optimization}}
\author{Peter Cotton}
\date{Draft, September 2026}

\begin{document}
\maketitle

\begin{abstract}
\noindent Derivative-free optimizers spend most of their budget on
one-dimensional line searches. We use the terminal-placement rule of
the companion theory paper as a line-search inner loop: with its scale
and correlation length estimated online, it flees, probes, or commits
in one to three evaluations. On the HumpDay physics benchmark it is a
competitive inner loop on rough and high-dimensional objectives at
tight budgets, where its early exits buy more search directions than a
classical line search. A projection argument then lifts the same rule
off the line: after two observations under any radial Gaussian-process
prior, the entire next-point decision lives in a plane, in any ambient
dimension. A fixed off-line placement does not beat a line search
across the catalog, but choosing it by expected improvement under a
fitted Gaussian process does, and increasingly so with dimension: the
value of the acquired point, not its position, is what turns the exact
reduction into a win.
\end{abstract}

\section{Line searches are where the budget goes}

Derivative-free optimizers reduce a high-dimensional problem to a
sequence of one-dimensional subproblems: pick a direction, place a few
evaluations along it, return a point. Classical inner loops, golden
section and Brent, assume smooth unimodal slices and spend eight to ten
evaluations resolving each one. On the rough slices of real physics
objectives that assumption fails, and on a tight global budget the cost
of each line is what binds.

The companion theory paper solves a different one-dimensional problem
exactly: place three evaluations on an exponentiated Ornstein--Uhlenbeck
slice and be paid the value at the final point. Its policy flees a
below-median start, probes between two good values, and stays above a
ceiling, all in closed form. Read as an inner loop it costs one to
three evaluations per line, not eight to ten. This paper asks what that
buys in an optimizer, and how far off the line the same geometry
reaches.

Two findings. As an inner loop the rule is competitive on rough and
high-dimensional objectives when the budget is tight, for a reason that
is about evaluation economy more than about the placement theorem. And
the placement geometry is not confined to a line: after two
observations under a radial prior the next-point decision collapses to
a plane, exactly, in any ambient dimension. A fixed off-line placement
does not pay, but one aimed by expected improvement under a fitted
Gaussian process, by the value rather than the information of the
acquired point, beats a line search on three quarters of high-dimensional
tasks, which is where the reduction should matter.

\section{The terminal-placement inner loop}

We implement the three-evaluation rule as an inner loop, \texttt{grass3}:
the exact final placement and opening step of the theory paper, with the
mean, scale, and correlation length estimated online from the
optimizer's own history. It flees or stays after one or two evaluations
when the rule dictates and places the third otherwise. It is a heuristic
adaptation, not the exact policy: it interpolates the opening-correlation
table, extrapolates it above the tabulated range, forces flight after
three stalled lines, and clips placements to the unit cube. The last is
a genuine approximation: choosing an unconstrained exterior winner and
then clipping can lose to the best feasible interior point, and we
flag it rather than hide it.

We benchmark \texttt{grass3} inside a common outer loop, iterated
random-direction search at budget 120, against four fixed inner loops:
golden section with two and with six evaluations, Brent, and random
placement, on matched directions and seeds.

\section{Benchmark}

The objectives are the application suite of the HumpDay package (Cotton
2021): small physics and engineering problems in which a point in the
unit cube sets a deterministic simulator's controls and the objective is
the outcome: a break in pool, pins felled in bowling, a trebuchet's
range, a wind-farm layout. We use the 22 dynamics-simulating problems;
their one-dimensional slices run from smooth to collision-cascade rough.

\begin{table}[t]
\centering
\small
\begin{tabular}{lccccc}
\toprule
Objective & $d$ & \texttt{grass3} & rank & best rival & seed wins \\
\midrule
plinko\_funnel & 7  & $-57.0$  & 1/5 & golden6 $-50.8$  & 13/16 \\
wind\_farm     & 16 & $-94.1$  & 1/5 & golden6 $-93.6$  & 16/24 \\
bowling        & 4  & $-105.0$ & 1/5 & golden2 $-104.0$ & 4/8   \\
free\_kick     & 4  & $-32.9$  & 5/5 & golden2 $-105.2$ & 3/24  \\
\bottomrule
\end{tabular}
\caption{\texttt{grass3} as an inner loop, median objective over seeds
at budget 120; lower is better. Comparators are golden section (two and
six evaluations) and Brent in the shared outer loop, and whole-cube
random search. Rank is among the five, averaging tied medians; seed
wins count seeds on which \texttt{grass3} beats the best rival. Across
all 22 objectives \texttt{grass3} has mean rank $2.95/5$ and is first or
second on eight (nine counting joint second places);
\texttt{free\_kick} is a representative loss.}
\label{tab:bench}
\end{table}

The rule wins on two kinds of objective. On collision and agent-chaos
objectives with rough but structured slices it is the best inner loop;
on a plinko objective it leads its strongest rival on 13 of 16 seeds
(Table~\ref{tab:bench}). It also wins on high-dimensional objectives even
when slices are smooth: on a sixteen-dimensional wind-farm layout it
leads golden6 on 16 of 24 seeds, and golden2 and Brent on 21 and 22.

The high-dimensional win is evaluation economy. Early exits of one or
two evaluations let the outer loop try several times more directions
than an inner loop spending eight to ten per line. The win is
budget-starved rather than general: it appears when the dimension times
a classical line's cost exceeds the whole budget, so a method with a
fixed number of directions per iteration cannot complete one pass.

The rule loses in two regimes. On landscapes too rough to carry
navigable structure, broad sampling wins and placement theory has
nothing to exploit; on needle landscapes, where the value hides in a
spike, golden section's fixed full-segment probes are the right move. It
is also not a drop-in for the line search inside an existing optimizer:
substituted into Powell, whose direction updates need genuine line
minimization, it loses to Brent.

Three qualifications keep the reading honest. Random search samples the
whole cube rather than sharing the outer loop, and the Brent baseline
spends one call reevaluating the known incumbent, so the economy margin
is somewhat overstated. An ablation is more sobering: omitting the third
placement altogether helps on twelve of the twenty-two objectives,
sometimes by a wide margin, so the benchmark speaks for evaluation
economy and online adaptation more than for the placement theorem in
particular. A control on exact exponentiated-OU fields, the model's own
landscape, loses to broad probing: terminal-placement wins do not
come from the objectives resembling an OU path.

As a standalone optimizer rather than an inner loop, a budget-truncated
two-shot form is mid-field against a panel of eleven standard methods,
best at the smallest budgets and overtaken by population methods as the
budget grows, the expected place for a rule built from a
few-evaluation terminal-placement model.

\section{From lines to planes}

The line is not special. Fix a Gaussian-process prior with constant mean
and radial covariance $K(\|x-y\|)$, and observe values at
$x_1,\dots,x_n$. Let $H$ be their affine hull and write a candidate as
$x=h+w$ with $h\in H$ and $w\perp H$. Then
\begin{equation}
\|x-x_i\|^{2}=\|h-x_i\|^{2}+\|w\|^{2},
\end{equation}
so the candidate's posterior depends only on its position within $H$ and
its distance from $H$. Every perpendicular direction is equivalent.

\begin{proposition}[Dimension reduction]
Under a radial Gaussian-process prior on an unrestricted domain, a
single-point acquisition can be optimized in at most $n$ dimensions
after $n$ observations, regardless of the ambient dimension.
\end{proposition}

With two observations the next-point decision lives in a plane, in any
ambient dimension, and needs no OU structure. The companion theory paper
computes that planar decision exactly for the exponential kernel: after
observing values $b,d$ at correlation $\rho$, the candidate's
correlations $(x,y)$ to the two observations range over the triangle
inequalities $\{xy\le\rho,\ x\ge\rho y,\ y\ge\rho x\}$, whose boundary is
the one-dimensional bracket and two exterior rays and whose interior
holds one extra candidate, the point that matches both observed values.

That extra candidate is worth a measurable amount. With
$b=d=\rho=\tfrac12$ the optimal third point forms an equilateral triangle
with the two observations, and its conditional expected terminal reward
exceeds the best collinear choice by about $2.9\%$. The gain is a planar
effect: it is unavailable to any line search.

\section{An adaptive planar search}

The reduction suggests a move. Evaluate a trial along a chosen direction;
use its observed value to choose the next point within a plane containing
that line; under the isotropic model, take any perpendicular direction.
The plane costs one extra coordinate and, by the dimension-reduction
proposition, captures the whole next-point decision no matter how large
the ambient space.

One adjustment is necessary, and the opening constant closes. The
theory paper rewards the terminal recommendation, whereas an optimizer
keeps the best evaluated point, so its opening maximizes
$\E[\max(e^{b},e^{X})]$ over the probe correlation $\rho$, with
$X\sim N(b\rho,1-\rho^{2})$. The geometric reduction is unaffected, a fact about the
posterior and not the objective. But the opening coefficient is not.

\begin{proposition}[Expected-improvement opening]
For the best-evaluated objective the optimal opening correlation
satisfies $\rho^{*}(b)=\kappa_{\mathrm{EI}}\,b+O(b^{3})$ as $b\downarrow0$,
with
\begin{equation}
\kappa_{\mathrm{EI}}=\frac{\E[e^{Z};Z>0]}{\E[Z\,e^{Z};Z>0]}
=\frac{\Phi(1)}{\Phi(1)+\phi(1)}=0.776639\ldots,
\qquad Z\sim N(0,1).
\end{equation}
\end{proposition}

\noindent The first-order condition
$\partial_\rho\E[\max(e^{b},e^{X})]
=\E\!\left[e^{X}\mathbf{1}\{X>b\}\,(b-\rho X'/\sqrt{1-\rho^{2}})\right]=0$,
with $X=b\rho+\sqrt{1-\rho^{2}}X'$ and $X'\sim N(0,1)$, gives at
$b\downarrow0$ the ratio
$\E[e^{Z};Z>0]=\kappa_{\mathrm{EI}}\,\E[Z e^{Z};Z>0]$, and
$\E[e^{Z};Z>0]=e^{1/2}\Phi(1)$, $\E[Z e^{Z};Z>0]=e^{1/2}(\phi(1)+\Phi(1))$.
So the optimizer opens \emph{wider} than the terminal rule,
$\kappa_{\mathrm{EI}}=0.777$ against the theory paper's
$\kappa_0=0.540$, because a fresh high draw is worth more when the
best point is kept rather than a mean recommendation. And it is
elementary, where $\kappa_0$ is only characterized by a transcendental
equation. A full planar-search implementation is the remaining
empirical step.

\section{The fixed off-line placement}

A budget-matched experiment isolates the planar move: two evaluations
per step along a shared random direction, differing only in where the
second goes. The \emph{line} keeps it on the line, extrapolating past
the trial when it improved and reflecting when it did not; the
\emph{plane} puts it at a fixed equilateral apex off the line. Start
point, direction stream, and adaptive step are identical. Across the
whole catalog of $51$ analytic objectives at $d\in\{2,5,10\}$ and $69$
physics simulators at native dimensions, $222$ tasks in all, the fixed
off-line move wins on only $0.26$ of seeds, below one half at every
dimension ($0.30$ at $d\le3$, $0.21$ at $d\ge8$), and a
momentum-estimated transverse direction is no better.

The fixed move ignores information. The reduction says the decision
lives in the plane, but a fixed apex in a random plane sends the one
off-line probe where the objective's structure almost never is; the
reduction is a fact about the posterior, not about any particular
objective.

\section{Aiming the plane by expected improvement}

Choose the off-line point by an acquisition instead of geometry. Fit a
Gaussian process to the evaluated history, and among candidate points
in the plane take the one a chosen acquisition ranks highest; the
control is the same GP restricted to the line, identical guidance and
budget, so the difference isolates the off-line freedom. Two
acquisitions separate cleanly (Table~\ref{tab:planargp}). Expected
improvement, which favors points likely to beat the incumbent, wins
increasingly with dimension. Pure information gain, the
posterior-variance criterion of entropy search (Hennig and Schuler
2012; Wang and Jegelka 2017) that favors the most uncertain point
regardless of value, does not.

\begin{table}[t]
\centering
\small
\begin{tabular}{lcc}
\toprule
dimension & expected improvement & information gain \\
\midrule
$d=2$--$3$ & $0.42$ & $0.27$ \\
$d=4$--$7$ & $0.63$ & $0.37$ \\
$d\ge 8$   & $0.76$ & $0.29$ \\
\midrule
overall    & $0.62$ & $0.31$ \\
\bottomrule
\end{tabular}
\caption{Fraction of seeds on which the GP-aimed off-line placement
beats a GP restricted to the line, over the $223$-task catalog at ten
seeds each. Both use the same GP and budget; only the off-line
candidates and the acquisition differ. Expected improvement helps more
as the dimension grows (analytic $0.78$, simulators $0.71$ at $d\ge8$);
the pure information-gain criterion does not help at any dimension.}
\label{tab:planargp}
\end{table}

Value, not information, turns the reduction into a win. The off-line
freedom pays only when the one probe is aimed at a point likely to be
good; aiming it at the most uncertain point, pure exploration, is
no better than the fixed geometry. The geometry hands high-dimensional
search a plane, and expected improvement decides where in it to look;
the payoff grows in exactly the regime where a line search has the
least room. All of this is heuristic: a greedy one-step acquisition
on a fitted Gaussian process, not the exact policy of the theory paper.
The only exact object it uses is the dimension reduction itself.

\section{Related work}

LineBO (Kirschner et al.\ 2019) optimizes through adaptively selected
one-dimensional subspaces; the reduction here identifies when one extra
transverse coordinate carries measurable value, sharpening the choice of
subspace to a plane after two observations. The Stochastic Three Points
method (Bergou et al.\ 2020) compares an incumbent with two opposite
trials; the planar candidate suggests an adaptive triangle variant whose
third location is chosen after seeing the second value. Knowledge
gradient (Frazier, Powell and Dayanik 2009) supplies the
measure-then-choose framework; the contribution here is explicit geometry
and the resulting planar move within it.

\section{Scope}

The defensible claim for high-dimensional derivative-free optimization is
an exact simplification of the next decision under a radial prior:
whatever the ambient dimension, two observations reduce it to a plane.
Learning useful directions over a long run remains hard, and constraints
or unknown anisotropy break the isotropy the reduction uses. The
immediate application is therefore small evaluation budgets and local
searches, where a single well-placed plane is most of what a step can
afford.

\begin{thebibliography}{9}

\bibitem{cottongrass2026}
P.~Cotton.
When the grass is greener: three-shot search on exponentiated Gaussian
landscapes.
Working paper, 2026.
\url{https://github.com/microprediction/browniansearch}.

\bibitem{humpday2021}
P.~Cotton.
HumpDay: pure Python derivative-free optimization and its application
suite.
\url{https://github.com/microprediction/humpday}, 2021.

\bibitem{kirschner2019}
J.~Kirschner, M.~Mutn\'y, N.~Hiller, R.~Ischebeck, and A.~Krause.
Adaptive and safe Bayesian optimization in high dimensions via
one-dimensional subspaces.
In \emph{Proceedings of the 36th International Conference on Machine
Learning}, 2019.

\bibitem{bergou2020}
E.~H. Bergou, E.~Gorbunov, and P.~Richt\'arik.
Stochastic three points method for unconstrained smooth minimization.
\emph{SIAM Journal on Optimization}, 30(4):2726--2749, 2020.

\bibitem{frazier2009}
P.~I. Frazier, W.~B. Powell, and S.~Dayanik.
The knowledge-gradient policy for correlated normal beliefs.
\emph{INFORMS Journal on Computing}, 21(4):599--613, 2009.

\bibitem{rasmussen2006}
C.~E. Rasmussen and C.~K.~I. Williams.
\emph{Gaussian Processes for Machine Learning}.
MIT Press, 2006.

\bibitem{hennig2012}
P.~Hennig and C.~J. Schuler.
Entropy search for information-efficient global optimization.
\emph{Journal of Machine Learning Research}, 13:1809--1837, 2012.

\bibitem{wang2017}
Z.~Wang and S.~Jegelka.
Max-value entropy search for efficient Bayesian optimization.
In \emph{Proceedings of the 34th International Conference on Machine
Learning}, 2017.

\end{thebibliography}

\end{document}
